% Generated draft: requires mathematical and source review.
\documentclass[11pt]{article}
\usepackage[margin=25mm]{geometry}
\usepackage{fontspec}
\setmainfont{Latin Modern Roman}
\usepackage{amsmath,amssymb,mathtools,unicode-math}
\setmathfont{Latin Modern Math}
\usepackage{enumitem,needspace,xurl,hyperref}
\setlength{\parindent}{0pt}
\setlength{\parskip}{5pt}
\setlength{\emergencystretch}{3em}
\newcommand{\N}{\mathbb{N}}
\newcommand{\Z}{\mathbb{Z}}
\newcommand{\Q}{\mathbb{Q}}
\newcommand{\R}{\mathbb{R}}
\newcommand{\C}{\mathbb{C}}
\newcommand{\F}{\mathbb{F}}
\newcommand{\A}{\mathbb{A}}
\newcommand{\PP}{\mathbb{P}}
\newcommand{\E}{\mathbb{E}}
\newcommand{\eps}{\varepsilon}
\newcommand{\dd}{\,\mathrm{d}}
\begin{document}
\section*{UnsolvedMath ID: 30004989}
\textbf{Source problem code:} catalogue entry
\subsection{Definitions and mathematical statement}
Let \(\mathfrak c=2^{\aleph_0}\) denote the cardinality of the continuum, and let \(\mathrm{CH}\) denote the assertion \(\mathfrak c=\aleph_1\). For an equivalence relation \(E\) on a set \(M\), write \([x]_E=\{y\in M:xEy\}\) and \(M/E=\{[x]_E:x\in M\}\). A quotient structure additionally carries the operations and relations induced by its presentation. The supplied record does not specify the signatures, presentation spaces, equivalence relations, or admissibility hypotheses intended by the phrase ``Borel quotient structures.''

For quotient structures in a common signature, let \(\operatorname{Aut}(M/E)\) denote the group of structure automorphisms and let \(\operatorname{Iso}(M/E,N/F)\) denote the set of structure isomorphisms. If topologies on the presentation spaces \(M\) and \(N\) have been specified, an isomorphism \(\Phi:M/E\to N/F\) has the lifted graph
\[
\Gamma_\Phi=\{(x,y)\in M\times N:\Phi([x]_E)=[y]_F\}.
\]
A customary interpretation of ``topologically trivial'' is that \(\Gamma_\Phi\) is Borel in the product of the presentation topologies. The supplied record does not define this term or provide those topologies, so this interpretation requires confirmation and is not adopted as an established source definition here. Write \(T_{E,F}(\Phi)\) for the intended topological-triviality property, pending that confirmation.

\textbf{Conjecture 3, as supplied, with unresolved scope.} For the intended admissible Borel quotient structures \(M/E\) and \(N/F\), the proposed assertions are
\[
\mathrm{CH}\ \Longrightarrow\ \bigl|\operatorname{Aut}(M/E)\bigr|=2^{\mathfrak c},
\]
and
\[
\text{the intended forcing axioms}\ \Longrightarrow\
\bigl(\forall\Phi\in\operatorname{Iso}(M/E,N/F)\bigr)\,T_{E,F}(\Phi).
\]
The source also states the following consequence of the first assertion:
\[
\mathrm{CH}\ \Longrightarrow\
\left|\{\alpha\in\operatorname{Aut}(M/E):\neg T_{E,E}(\alpha)\}\right|=2^{\mathfrak c}.
\]
The presentation assumptions needed to justify this consequence, including a suitable bound on the number of topologically trivial automorphisms, are not specified in the supplied record.

These are schematic assertions: the precise class of admissible quotients, its quantifiers and nondegeneracy hypotheses, the definition of topological triviality, and the particular forcing axioms remain unresolved. In particular, the record does not support an unqualified universal assertion over all Borel quotient structures.

The passage says that the work mainly concerns quotients of Boolean algebras and certain quotients of \(C^*\)-algebras called corona \(C^*\)-algebras, described there as noncommutative generalizations of \v{C}ech--Stone remainders of locally compact noncompact spaces. It reports instances of the conjecture for large classes of these structures, with the work contained in a survey labeled [1] in the passage. These contextual statements do not identify the exact hypotheses of those instances.
\subsection{Short English statement}
The supplied conjectural pattern predicts \(2^{\mathfrak c}\) automorphisms, including that many topologically nontrivial ones, under CH, and topological triviality of every isomorphism under unspecified forcing axioms. The intended quotient classes, presentation hypotheses, triviality definition, and forcing axioms must be recovered before this becomes a fully precise mathematical statement. The supplied status is ``partially solved''; this draft does not independently verify that assessment.
\subsection{Sources}
\begin{itemize}
\item Local UnsolvedMath registry: supplied complete source record, permanent ID \(30004989\), source code \texttt{catalogue entry}; registry paths \texttt{data/unsolvedmath/problems.json} and \texttt{catalog/problems.json}. Title: ``Set-Theoretic Dichotomy for Borel-Quotient Automorphisms.'' No individual proposer is identified in the record.
\item Supplied report citation: \emph{Set Theory (2023)}, Oberwolfach Reports, pp.\ 79--106, DOI \url{https://doi.org/10.4171/owr/2022/2}. The extracted passage is labeled Conjecture 3. The record gives proposed year 2022; the relationship between this year and the supplied report title's year is not resolved here.
\item Supplied literature reference: Farah--Ghasemi--Vaccaro--Vignati, \emph{Corona Rigidity}, \url{https://arxiv.org/abs/2201.11618}.
\item Supplied literature reference: ``2025 Bulletin of Symbolic Logic version,'' \url{https://doi.org/10.1017/bsl.2025.10084}.
\item Literature assessment dated exactly \texttt{2026-08-22}: the supplied curation labels the problem ``partially solved'' and reports that the CH/forcing-axiom pattern is verified for broad families of Boolean, \(C^*\), and corona quotients, while the full Borel-quotient conjecture remains a survey-level open problem. This is supplied dated curation, not independently verified literature assessment. The background explicitly requests independent verification of time-sensitive status and cited claims.
\item The record labels its statement \texttt{corrected\_verified} and reports that the official Oberwolfach report was checked directly. This is a supplied verification claim, not a source visit or verification performed for this draft. The passage refers to a survey as [1], but its numbered bibliography entry is not reproduced in the supplied record.
\end{itemize}
\end{document}
